\subsection{Predicted Amplitude and Redshift Dependence}
  \label{subsec:hubble_predictions}

  A key feature of the present framework is that the magnitude of the
  locally inferred deviation from the global expansion rate is
  constrained by the saturation scale $a_\star$, which is fixed
  independently from galaxy rotation curves.

  For representative underdensities characteristic of large
  cosmic voids on scales of a few hundred megaparsecs,
  and for $g_{\mathrm{void}} \ll a_\star$,
  the framework predicts a fractional enhancement
  \begin{equation}
    \frac{\Delta H}{H}
    \;\equiv\;
    \frac{H_{\mathrm{loc}} - H_{\mathrm{global}}}
    {H_{\mathrm{global}}}
    \;\sim\;
    \text{a few percent},
    \label{eq:deltaH}
  \end{equation}
  with the precise value controlled by the empirically
  determined scale $a_\star$ and by independently measured
  void density contrasts.

  The effect is intrinsically redshift dependent.
  At higher redshift, large-scale structure is less developed,
  density contrasts are reduced, and line-of-sight averaging
  over multiple environments becomes increasingly efficient.
  Consequently, the locally inferred expansion rate is expected
  to converge progressively toward $H_{\mathrm{global}}$
  with increasing redshift.

  The framework therefore predicts a smooth decrease of
  $\Delta H/H$ as a function of survey depth.
  A persistent offset at high redshift, inconsistent with
  environmental averaging, would falsify the mechanism.

  Conversely, deviations substantially larger than the
  percent-to-few-percent range, or the complete absence of any
  measurable environment-dependent trend once systematics are
  controlled, would directly challenge the proposed interpretation.
